\documentclass[11pt,a4paper]{article}
\usepackage[utf8]{inputenc}
\usepackage[T1]{fontenc}
\usepackage[french]{babel}
\usepackage{amsmath}
\usepackage{amssymb}
\usepackage{amsthm}
\usepackage{mathtools}
\usepackage{geometry}
\usepackage{enumitem}
\usepackage{xcolor}
\usepackage{titlesec}
\usepackage{fancyhdr}
\usepackage{booktabs}
\usepackage{array}
\usepackage{tcolorbox}

% Configuration de la page
\geometry{margin=2.5cm}
\pagestyle{fancy}
\fancyhf{}
\fancyhead[L]{\leftmark}
\fancyhead[R]{Exercice P-value}
\fancyfoot[C]{\thepage}

% Configuration des titres
\titleformat{\section}
{\Large\bfseries\color{blue!70!black}}
{}
{0em}
{}[\titlerule]

\titleformat{\subsection}
{\large\bfseries\color{blue!50!black}}
{}
{0em}
{}

\titleformat{\subsubsection}
{\normalsize\bfseries\color{blue!30!black}}
{}
{0em}
{}

% Boîtes colorées
\newtcolorbox{astucebox}{
    colback=green!5!white,
    colframe=green!75!black,
    title=\textbf{Astuce},
    fonttitle=\bfseries
}

\newtcolorbox{rappelbox}{
    colback=blue!5!white,
    colframe=blue!75!black,
    title=\textbf{Rappel},
    fonttitle=\bfseries
}

\newtcolorbox{attentionbox}{
    colback=orange!5!white,
    colframe=orange!75!black,
    title=\textbf{Attention},
    fonttitle=\bfseries
}

% Commandes personnalisées
\newcommand{\R}{\mathbb{R}}
\newcommand{\Var}{\text{Var}}
\newcommand{\E}{\mathbb{E}}
\newcommand{\N}{\mathcal{N}}
\newcommand{\student}{t}

% Métadonnées
\title{Exercice Complet - P-value\\
\large Interprétation et Calcul}
\author{AmineGR03}
\date{\today}

\begin{document}

\maketitle

\tableofcontents
\newpage

\section{Introduction}

Ce document présente un exercice complet sur la p-value, avec une explication détaillée de son calcul et de son interprétation dans le contexte d'un test d'hypothèse.

\newpage

\section{Exercice : Test sur la Moyenne avec P-value}

\subsection{Énoncé}

Un fabricant de batteries affirme que ses batteries ont une durée de vie moyenne d'au moins 120 heures. Un contrôleur qualité soupçonne que la durée de vie moyenne est inférieure à cette valeur. Il teste un échantillon de 25 batteries et obtient les résultats suivants :
\begin{itemize}
    \item Durée de vie moyenne observée : $\overline{X} = 115$ heures
    \item Écart-type observé : $s = 10$ heures
    \item Taille de l'échantillon : $n = 25$
\end{itemize}

On suppose que la durée de vie suit une loi normale.

\begin{enumerate}
    \item Formuler les hypothèses du test.
    \item Calculer la statistique de test.
    \item Déterminer la p-value.
    \item Interpréter la p-value.
    \item Conclure au seuil $\alpha = 0.05$.
    \item Que se passerait-il si on choisissait $\alpha = 0.01$ ?
    \item Donner un intervalle de confiance à 95\% pour la durée de vie moyenne.
\end{enumerate}

\subsection{Rappels Théoriques}

\begin{rappelbox}
\textbf{P-value (valeur-p) :}

La p-value est la probabilité d'observer une statistique de test au moins aussi extrême que celle observée, sous l'hypothèse que $H_0$ est vraie.

\textbf{Interprétation :}
\begin{itemize}
    \item Si $p < \alpha$ : on rejette $H_0$ (résultat significatif)
    \item Si $p \geq \alpha$ : on ne rejette pas $H_0$ (pas de preuve contre $H_0$)
\end{itemize}

\textbf{Test sur la moyenne (variance inconnue) :}

Si $\sigma$ est inconnu et $n$ petit, on utilise la loi de Student :
$$t = \frac{\overline{X} - \mu_0}{s/\sqrt{n}} \sim t_{n-1}$$

\textbf{Test unilatéral à gauche :}
\begin{itemize}
    \item $H_0 : \mu \geq \mu_0$ contre $H_1 : \mu < \mu_0$
    \item P-value : $p = P(T_{n-1} \leq t)$
\end{itemize}
\end{rappelbox}

\subsection{Solution Détaillée}

\begin{enumerate}
    \item \textbf{Formulation des hypothèses :}
    
    Le fabricant affirme "au moins 120 heures", donc :
    \begin{align}
    H_0 &: \mu \geq 120 \quad \text{(la durée de vie moyenne est d'au moins 120 heures)} \\
    H_1 &: \mu < 120 \quad \text{(la durée de vie moyenne est inférieure à 120 heures)}
    \end{align}
    
    C'est un \textbf{test unilatéral à gauche}.
    
    \item \textbf{Calcul de la statistique de test :}
    
    Comme $\sigma$ est inconnu et $n = 25$ (petit échantillon), on utilise la loi de Student avec $\nu = n - 1 = 24$ degrés de liberté.
    
    $$t = \frac{\overline{X} - \mu_0}{s/\sqrt{n}} = \frac{115 - 120}{10/\sqrt{25}} = \frac{-5}{10/5} = \frac{-5}{2} = -2.5$$
    
    \item \textbf{Calcul de la p-value :}
    
    Pour un test unilatéral à gauche, la p-value est :
    $$p = P(T_{24} \leq -2.5) = P(T_{24} \geq 2.5)$$
    
    (Par symétrie de la loi de Student)
    
    En consultant la table de Student avec $\nu = 24$ degrés de liberté :
    \begin{itemize}
        \item $P(T_{24} \geq 2.492) = 0.01$
        \item $P(T_{24} \geq 2.797) = 0.005$
    \end{itemize}
    
    Par interpolation linéaire entre $t = 2.492$ et $t = 2.797$ :
    \begin{align}
    p &\approx 0.01 - \frac{0.01 - 0.005}{2.797 - 2.492} \times (2.5 - 2.492) \\
    &\approx 0.01 - \frac{0.005}{0.305} \times 0.008 \\
    &\approx 0.01 - 0.00013 \\
    &\approx 0.00987
    \end{align}
    
    Donc $p \approx 0.010$ (environ 1\%)
    
    \item \textbf{Interprétation de la p-value :}
    
    La p-value de $0.010$ signifie que :
    \begin{itemize}
        \item Si $H_0$ était vraie (si la durée de vie moyenne était vraiment $\geq 120$ heures)
        \item La probabilité d'observer une moyenne d'échantillon aussi faible (ou plus faible) que 115 heures
        \item Est d'environ 1\%
    \end{itemize}
    
    C'est une probabilité \textbf{faible}, ce qui suggère que les données observées sont peu compatibles avec $H_0$.
    
    \item \textbf{Conclusion au seuil $\alpha = 0.05$ :}
    
    Puisque $p = 0.010 < \alpha = 0.05$, on \textbf{rejette $H_0$} au seuil de 5\%.
    
    On conclut qu'il existe des \textbf{preuves statistiques significatives} que la durée de vie moyenne des batteries est \textbf{inférieure à 120 heures}.
    
    \item \textbf{Conclusion au seuil $\alpha = 0.01$ :}
    
    Puisque $p = 0.010 = \alpha = 0.01$, on est à la limite. En général, on considère que si $p = \alpha$, on rejette $H_0$.
    
    Cependant, avec $p \approx 0.010$, on est très proche de la frontière. On peut dire qu'on \textbf{rejette $H_0$} au seuil de 1\%, mais de justesse.
    
    \textbf{Note :} Si on avait $p = 0.011 > 0.01$, on ne rejetterait pas $H_0$ au seuil de 1\%.
    
    \item \textbf{Intervalle de confiance à 95\% :}
    
    Pour un niveau de confiance de 95\%, on a $\alpha = 0.05$, donc $\alpha/2 = 0.025$.
    
    Avec $\nu = 24$ degrés de liberté, on consulte la table : $t_{0.025, 24} \approx 2.064$
    
    L'intervalle de confiance est :
    $$IC_{95\%} = \left[\overline{X} - t_{0.025, 24} \frac{s}{\sqrt{n}}, \overline{X} + t_{0.025, 24} \frac{s}{\sqrt{n}}\right]$$
    
    $$IC_{95\%} = \left[115 - 2.064 \times \frac{10}{\sqrt{25}}, 115 + 2.064 \times \frac{10}{\sqrt{25}}\right]$$
    
    $$IC_{95\%} = \left[115 - 2.064 \times 2, 115 + 2.064 \times 2\right]$$
    
    $$IC_{95\%} = \left[115 - 4.128, 115 + 4.128\right] = [110.87, 119.13]$$
    
    \textbf{Interprétation :} On est confiant à 95\% que la durée de vie moyenne des batteries se situe entre 110.87 et 119.13 heures.
    
    On remarque que $\mu_0 = 120$ n'est \textbf{pas} dans l'intervalle de confiance, ce qui confirme le rejet de $H_0$.
\end{enumerate}

\subsection{Tableau Récapitulatif}

\begin{table}[h]
\centering
\small
\begin{tabular}{|l|c|}
\hline
\textbf{Élément} & \textbf{Valeur} \\
\hline
Hypothèse nulle $H_0$ & $\mu \geq 120$ \\
\hline
Hypothèse alternative $H_1$ & $\mu < 120$ \\
\hline
Moyenne observée $\overline{X}$ & 115 heures \\
\hline
Écart-type observé $s$ & 10 heures \\
\hline
Taille échantillon $n$ & 25 \\
\hline
Statistique de test $t$ & $-2.5$ \\
\hline
Degrés de liberté $\nu$ & 24 \\
\hline
P-value $p$ & $\approx 0.010$ (1\%) \\
\hline
Seuil $\alpha$ & 0.05 (5\%) \\
\hline
Décision ($\alpha = 0.05$) & Rejeter $H_0$ \\
\hline
Décision ($\alpha = 0.01$) & Rejeter $H_0$ (de justesse) \\
\hline
\end{tabular}
\caption{Résumé du test}
\end{table}

\begin{attentionbox}
\textbf{Erreurs fréquentes à éviter :}
\begin{itemize}
    \item La p-value n'est \textbf{PAS} la probabilité que $H_0$ soit vraie
    \item La p-value est la probabilité d'observer des données aussi extrêmes \textbf{sous l'hypothèse que $H_0$ est vraie}
    \item Une p-value faible ne signifie pas que l'effet est important, seulement qu'il est statistiquement significatif
    \item Ne pas confondre significativité statistique et significativité pratique
\end{itemize}
\end{attentionbox}

\begin{astucebox}
\textbf{Interprétation pratique de la p-value :}
\begin{itemize}
    \item $p < 0.01$ : Preuve très forte contre $H_0$ (très significatif)
    \item $0.01 \leq p < 0.05$ : Preuve forte contre $H_0$ (significatif)
    \item $0.05 \leq p < 0.10$ : Preuve modérée contre $H_0$ (marginalement significatif)
    \item $p \geq 0.10$ : Pas de preuve contre $H_0$ (non significatif)
\end{itemize}
\end{astucebox}

\newpage

\section{Résumé des Étapes}

\begin{enumerate}
    \item \textbf{Formuler les hypothèses} : $H_0$ vs $H_1$
    \item \textbf{Choisir la statistique de test} : selon que $\sigma$ est connu ou non
    \item \textbf{Calculer la statistique observée} : $t$ ou $z$
    \item \textbf{Calculer la p-value} : selon le type de test (unilatéral ou bilatéral)
    \item \textbf{Interpréter la p-value} : probabilité sous $H_0$
    \item \textbf{Comparer à $\alpha$} : $p < \alpha$ $\Rightarrow$ rejeter $H_0$
    \item \textbf{Conclure} : dans le contexte du problème
\end{enumerate}

\vspace{2cm}

\begin{center}
\textbf{Fin du document}
\end{center}

\end{document}

